% ==============================================================================
% فصل دوم: مفاهیم پایه و پیشینه پژوهش
% ==============================================================================

\chapter{مفاهیم پایه و پیشینه پژوهش}
\label{ch:background}

\section{مقدمه}
ارزیابی و تضمین رفتارهای زمانی در سامانه‌های تعبیه‌شده بی‌درنگ، نیازمند فهم هم‌زمان مبانی ریاضی نظریه زمان‌بندی، ویژگی‌های ریزمعماری سخت‌افزار و استانداردهای صنعتی حاکم بر ردگیری دستورات است. در این فصل، ابتدا فرمول‌بندی ریاضی و مفاهیم بنیادین «بیشینه زمان اجرا» (\lr{WCET}) و طبقه‌بندی رویکردهای استنتاج آن مورد بررسی قرار می‌گیرد. سپس معماری عیب‌یابی و پایش بی‌درنگ \lr{ARM CoreSight} به عنوان استاندارد پیشرو در پردازنده‌های نوین تشریح شده و نقش هر یک از اجزای داخلی آن در زنجیره انتقال داده تبیین می‌گردد. در پایان، رویکرد این رساله در چارچوب پلتفرم‌های تجاری آزمون محصول و پیشینه پژوهش‌های بین‌المللی به صورت تطبیقی مقایسه خواهد شد.

\section{مفهوم فرمالیزه‌شده بیشینه زمان اجرا (\lr{WCET})}
در یک سامانه تعبیه‌شده بی‌درنگ، هر وظیفه محاسباتی (\lr{Task}) باید پیش از رسیدن به ضرب‌الاجل زمانی خود (\lr{Deadline}) خاتمه یابد. زمان اجرای یک وظیفه، کمیت ثابتی نیست بلکه تابعی از ورودی‌های متغیر محیطی و حالات داخلی سخت‌افزار است. 

به صورت ریاضی، اگر فضای ورودی‌های معتبر برنامه را با $\mathcal{P}$ و فضای تمامی حالات ممکن سخت‌افزار در آغاز اجرای وظیفه (شامل وضعیت خط‌لوله، حافظه نهان و صف‌های انتقال) را با $\mathcal{S}_0$ نمایش دهیم، زمان واقعی اجرای برنامه به صورت یک تابع غیرخطی به فرم زیر بیان می‌شود:
\begin{equation}
T = f(p, s_0) \quad \text{\rl{که در آن}} \quad p \in \mathcal{P}, \; s_0 \in \mathcal{S}_0
\end{equation}

بر این اساس، بیشینه زمان اجرا (\lr{WCET}) عبارت است از سوپرمم (بیشترین مقدار) این تابع بر روی تمامی ترکیبات ممکن ورودی‌ها و حالات اولیه سخت‌افزار:
\begin{equation}
\mathrm{WCET} \triangleq \max_{p \in \mathcal{P}, \; s_0 \in \mathcal{S}_0} f(p, s_0)
\end{equation}

\begin{figure}[ht]
\centering
\begin{latin}
\begin{tikzpicture}[xscale=1.18, yscale=1.05, every node/.append style={font=\latinfont}]
    % Axes
    \draw[-latex, thick] (0,0) -- (9.6,0) node[right=3pt] {\small \rl{زمان اجرا} (\lr{Time})};
    \draw[-latex, thick] (0,0) -- (0,4.4) node[above=2pt] {\small \rl{چگالی احتمال} $P(T)$};
    
    % Distribution curve
    \draw[very thick, blue!80!black, fill=blue!10, smooth, domain=0.8:6.0] plot (\x, {3.2*exp(-0.9*(\x-3.0)^2)});
    
    % Points on x axis
    \coordinate (BCET) at (1.4,0);
    \coordinate (ACET) at (3.0,0);
    \coordinate (MAX_OBS) at (4.8,0);
    \coordinate (WCET_REAL) at (6.4,0);
    \coordinate (WCET_EST) at (8.2,0);
    
    % Vertical lines
    \draw[dashed, gray] (1.4,0) -- (1.4,1.2);
    \draw[dashed, blue] (3.0,0) -- (3.0,3.2);
    \draw[thick, orange!90!black] (4.8,0) -- (4.8,2.2) node[above=3pt, font=\scriptsize\latinfont, text=orange!90!black, align=center] {\rl{بیشینه مشاهده‌شده}\\\lr{(Max Observed)}};
    \draw[thick, red!80!black] (6.4,0) -- (6.4,3.0) node[above=3pt, font=\scriptsize\latinfont, text=red!80!black, align=center] {\rl{بیشینه واقعی}\\\lr{(True WCET)}};
    \draw[thick, green!60!black] (8.2,0) -- (8.2,3.8) node[above=3pt, font=\scriptsize\latinfont, text=green!50!black, align=center] {\rl{کران امن برآوردشده}\\\lr{(Safe WCET Bound)}};
    
    % Labels below axis
    \node[below=4pt, font=\small\latinfont] at (BCET) {BCET};
    \node[below=4pt, font=\small\latinfont] at (ACET) {ACET};
    \node[below=4pt, font=\small\latinfont] at (MAX_OBS) {$T_{\mathrm{max\_obs}}$};
    \node[below=4pt, font=\small\latinfont] at (WCET_REAL) {$\mathrm{WCET}_{\mathrm{real}}$};
    \node[below=4pt, font=\small\latinfont] at (WCET_EST) {$\mathrm{WCET}_{\mathrm{bound}}$};
    
    % Risk areas - cleanly separated by y-level and color
    \draw[latex-latex, red!90!black, very thick] (4.8,0.7) -- node[above=4pt, font=\scriptsize\latinfont, text=red!90!black] {\rl{ناحیه ریسک ناامن}} (6.4,0.7);
    \draw[latex-latex, green!60!black, very thick] (6.4,1.6) -- node[above=4pt, font=\scriptsize\latinfont, text=green!60!black] {\rl{حاشیه اطمینان مهندسی}} (8.2,1.6);
\end{tikzpicture}
\end{latin}
\caption{فضای توزیع زمان اجرای برنامه و موقعیت بیشینه زمان اجرا نسبت به زمان‌های مشاهده‌شده و کران امن}
\label{fig:wcet_distribution}
\end{figure}

همان‌طور که در شکل \ref{fig:wcet_distribution} نشان داده شده است، در عمل یافتن دقیق نقطه $\mathrm{WCET}_{\mathrm{real}}$ به دلیل بی‌شمار بودن فضای حالات ورودی و سخت‌افزار بسیار دشوار است. بنابراین، مهندسی سامانه‌های بی‌درنگ به دنبال یافتن یک \textbf{کران بالای امن (\lr{Safe Upper Bound})} است به طوری که:
\begin{equation}
\mathrm{WCET}_{\mathrm{bound}} \ge \mathrm{WCET}_{\mathrm{real}}
\end{equation}

در صورتی که کران برآوردشده کمتر از مقدار واقعی باشد ($\mathrm{WCET}_{\mathrm{bound}} < \mathrm{WCET}_{\mathrm{real}}$)، با یک برآورد نامعتبر و ناامن (\lr{Unsafe Estimation}) روبه‌رو هستیم که می‌تواند موجب شکست سیستم در سناریوهای بحرانی شود. از طرف دیگر، اگر این کران بیش‌از‌حد محافظه‌کارانه و بزرگ تعیین شود (\lr{Over-estimation})، مهندس طراح مجبور به انتخاب پردازنده‌های گران‌تر و منابع سخت‌افزاری مازاد بر نیاز خواهد شد.

\section{طبقه‌بندی متدولوژی‌های استخراج \lr{WCET}}
رویکردهای تعیین بیشینه زمان اجرا در ادبیات علمی به سه شاخه بنیادین تقسیم می‌شوند \cite{wilhelm2008worst}:

\subsection{روش‌های تحلیل ایستا (\lr{Static WCET Analysis})}
در روش‌های ایستا، کد باینری کامپایل‌شده بدون اجرای واقعی بر روی سخت‌افزار، مورد پردازش و تحلیل صوری ریاضی قرار می‌گیرد. این متدولوژی شامل سه گام کلیدی است:
\begin{enumerate}
    \item \textbf{استخراج گراف جریان کنترل (\lr{Control Flow Graph - CFG}):} بازسازی ساختار شاخه‌ها، حلقه‌ها و توابع از کد ماشین.
    \item \textbf{مدل‌سازی سخت‌افزار (\lr{Hardware Modeling}):} پیش‌بینی رفتار چرخه‌ای پردازنده (تخمین رفتار خط‌لوله و نرخ اصابت کش با استفاده از تئوری تفسیر انتزاعی یا \lr{Abstract Interpretation}).
    \item \textbf{تکنیک شمارش ضمنی مسیرها (\lr{IPET - Implicit Path Enumeration}):} فرمول‌بندی مسئله یافتن طولانی‌ترین مسیر در قالب یک مسئله برنامه‌ریزی خطی عدد صحیح (\lr{Integer Linear Programming - ILP}). اگر $x_i$ بیانگر تعداد دفعات اجرای بلوک پایه $i$ و $c_i$ کران بالای زمان اجرای آن بلوک باشد، مسئله به فرم زیر بهینه‌سازی می‌شود:
    \begin{equation}
    \mathrm{WCET} = \max \sum_{i=1}^{N} c_i \cdot x_i
    \end{equation}
    با قید برقراری معادلات پیوستگی جریان ورودی و خروجی در گره‌های گراف جریان برنامه.
\end{enumerate}

\textbf{نقاط ضعف روش‌های ایستا:} با ظهور پردازنده‌های مدرن با هسته‌های پیچیده، مدل‌سازی رفتار چرخه به چرخه ریزمعماری به دلیل پدیده‌هایی نظیر عدم‌تعین حافظه‌های نهان اشتراکی، تداخل گذرگاه‌ها و اجرای خارج از نوبت (\lr{Out-of-Order Execution}) عملاً ناممکن گشته یا به تخمین‌هایی با خطای بالای ۵۰ تا ۱۰۰ درصد محافظه‌کاری منجر می‌شود.

\subsection{روش‌های مبتنی بر اندازه‌گیری تجربی (\lr{MBTA})}
در رویکرد \lr{MBTA}، برنامه با مجموعه‌ای از سناریوهای آزمون بر روی سخت‌افزار فیزیکی اجرا شده و زمان‌های پاسخ مستقیماً ثبت می‌گردند \cite{bernat2002wcet}. چالش‌های اصلی این متدولوژی عبارتند از:
\begin{itemize}
    \item \textbf{پوشش ورودی‌های بدترین شرایط (\lr{Worst-Case Input Space}):} تضمین اینکه آزمون‌ها توانسته‌اند مسیر منجر به بیشینه زمان را فعال سازند بسیار دشوار است (نقشی که در پروژه حاضر به سامانه تست خودکار آقای حسن آملی واگذار شده است).
    \item \textbf{اثر تهاجمی ابزار دقیق (\lr{Probe Effect}):} برای آگاهی از مسیر پیموده‌شده و زمان قطعات کد، نباید ساختار کد تغییر کند؛ معضلی که تنها با ردگیری سخت‌افزاری قابل حل است.
\end{itemize}

\subsection{روش‌های ترکیبی (\lr{Hybrid WCET Analysis})}
این روش‌ها از تلفیق نقاط قوت دو رویکرد پیشین پدید آمده‌اند \cite{betts2010hybrid}. در رویکرد ترکیبی، سنجش زمان اجرای قطعات پایه برنامه‌ریزی‌شده (\lr{Basic Blocks}) بر عهده ردگیری سخت‌افزاری بر روی پردازنده واقعی گذاشته می‌شود تا رفتار پیچیده خط‌لوله و کش مستقیماً از واقعیت استخراج گردد؛ سپس با تزریق این زمان‌های تجربی دقیق به عنوان ضرایب $c_i$ در مدل‌های ایستا (\lr{IPET})، بیشینه زمان اجرای کل ساختار برنامه به صورت ریاضی محاسبه می‌گردد.

\section{معماری عیب‌یابی و ردگیری \lr{ARM CoreSight}}
برای پاسخ به نیاز صنعت در زمینه پایش و صحه‌گذاری سامانه‌های بی‌درنگ بدون متوقف کردن هسته پردازنده یا تحمیل اثر پروب، شرکت \lr{ARM} معماری جامع \textbf{\lr{CoreSight}} را به عنوان استاندارد سخت‌افزاری درون‌تراشه‌ای پایه‌گذاری نمود \cite{armcoresight2013}. 

معماری \lr{CoreSight} یک زیرسیستم سخت‌افزاری کاملاً مجزا از هسته محاسباتی است که به صورت موازی در کنار هسته اجرا می‌شود. در تراشه‌های سری \lr{STM32H750} (مبتنی بر هسته \lr{Cortex-M7})، این زیرسیستم شامل زنجیره‌ای از واحدهای تخصصی است که در شکل \ref{fig:coresight_pipeline} نشان داده شده‌اند.

\begin{figure}[ht]
\centering
\begin{latin}
\begin{tikzpicture}[
    every node/.append style={font=\latinfont},
    node distance=2.5cm and 2.6cm,
    hw/.style={rectangle, draw=black!80, fill=gray!10, thick, rounded corners=5pt, text width=4.8cm, align=center, minimum height=1.75cm},
    trace/.style={rectangle, draw=blue!80, fill=blue!10, thick, rounded corners=5pt, text width=4.8cm, align=center, minimum height=1.75cm},
    ext/.style={rectangle, draw=orange!80, fill=orange!10, thick, rounded corners=5pt, text width=5.2cm, align=center, minimum height=1.75cm},
    arrow/.style={-latex, very thick}
]
    % On-chip Core & Trigger
    \node [hw] (core) {\textbf{Cortex-M7 Core}\\[0.15cm]\small Pipeline \& Execution};
    \node [hw, right=2.6cm of core] (dwt) {\textbf{DWT Unit}\\[0.15cm]\small Comparators (COMP0/1)};
    
    % Trace Generation
    \node [trace, below=2.5cm of core] (etm) {\textbf{ETMv4}\\[0.15cm]\small Trace Compression};
    
    % Transport
    \node [trace, below=2.5cm of etm] (cstf) {\textbf{CSTF (Trace Funnel)}\\[0.15cm]\small ATB Bus Multiplexing};
    \node [trace, right=2.6cm of cstf] (etf) {\textbf{ETF (Trace FIFO)}\\[0.15cm]\small Burst Data Buffer};
    \node [trace, below=2.5cm of cstf] (tpiu) {\textbf{TPIU}\\[0.15cm]\small Port Formatter (ATB \ensuremath{\to} Pins)};
    
    % Off-chip Pins & Analyzer
    \node [ext, below=2.7cm of tpiu] (pins) {\textbf{Physical Package Pins}\\[0.15cm]\small TRACECLK + TRACED[3:0]};
    \node [ext, right=2.8cm of pins] (dslogic) {\textbf{DSLogic U3Pro16}\\[0.15cm]\small \mbox{10x Oversampled} Logic Capture};
    \node [ext, above=2.7cm of dslogic] (host) {\textbf{Host Python / OpenCSD}\\[0.15cm]\small Snapshot Replay \& Decode};

    % Connections
    \draw [arrow] (core) -- node[left=4pt, font=\small] {Instruction Flow} (etm);
    \draw [arrow] (dwt.south) |- node[near end, above=4pt, font=\small] {Trigger Signals} (etm.east);
    \draw [arrow] (etm) -- node[left=4pt, font=\small] {ATB Stream} (cstf);
    \draw [arrow, latex-latex] (cstf) -- (etf);
    \draw [arrow] (cstf) -- (tpiu);
    \draw [arrow] (tpiu) -- node[left=4pt, font=\small] {4-bit Parallel} (pins);
    \draw [arrow] (pins) -- node[midway, above=4pt, font=\small, fill=white, inner sep=2pt] {Physical Probing} (dslogic);
    \draw [arrow] (dslogic) -- node[right=4pt, font=\small] {Raw CSV Samples} (host);
\end{tikzpicture}
\end{latin}
\caption{زنجیره کامل جریان سیگنال‌ها در معماری \lr{ARM CoreSight} از هسته پردازنده تا ابزار تحلیل}
\label{fig:coresight_pipeline}
\end{figure}

\subsection{اجزای زیرسیستم \lr{CoreSight} در \lr{STM32H750}}
\begin{enumerate}
    \item \textbf{واحد پورت دسترسی به عیب‌یابی (\lr{DAP - Debug Access Port}):} درگاه اصلی اتصال ابزارهای دیباگ خارجی از طریق رابط‌های دو سیمه \lr{SWD} یا چندسیمه \lr{JTAG} جهت دسترسی به ثبّات‌های کنترلی و فضای آدرس حافظه بدون مداخله هسته.
    \item \textbf{گذرگاه پیشرفته ردگیری (\lr{ATB - Advanced Trace Bus}):} پروتکل گذرگاه داده همگام و پرسرعت درون‌تراشه‌ای که بسته‌های ردگیری تولیدشده توسط ماژول‌های پایش را میان اجزای گوناگون منتقل می‌سازد.
    \item \textbf{واحد نگهبان داده و ردگیری (\lr{DWT - Data Watchpoint and Trace}):} ماژولی حیاتی که وظایف متعددی نظیر شمارش چرخه‌های زمانی (\lr{CYCCNT}) و مقایسه آدرس‌های حافظه را انجام می‌دهد. ثبّات‌های مقایسه‌گر \lr{DWT\_COMP0} و \lr{DWT\_COMP1} نقشی محوری در تعریف محدوده پنجره‌های پایش دارند؛ به طوری که با رسیدن آدرس اشاره‌گر دستورات به یک برچسب خاص، سیگنال شروع یا پایان ردگیری به \lr{ETM} صادر می‌گردد.
    \item \textbf{قیف ادغام ردگیری (\lr{CSTF - CoreSight Trace Funnel}):} ماژولی چندورودی-تک‌خروجی که وظیفه تلفیق جریان داده‌های ردگیری از منابع مختلف درون‌تراشه‌ای (نظیر هسته‌های چندگانه، \lr{ETM} و \lr{ITM}) به یک گذرگاه واحد \lr{ATB} را با شناسه مبدأ (\lr{Trace Source ID}) مشخص عهده‌دار است.
    \item \textbf{بافر میانی جریان داده (\lr{ETF - Embedded Trace FIFO}):} یک حافظه \lr{SRAM} حلقوی درون‌تراشه‌ای جهت جذب انفجارهای ناگهانی داده‌های \lr{Trace} (\lr{Data Bursts})؛ در زمان‌هایی که انشعاب‌های شرطی مکرر موجب افزایش نرخ تولید بسته شده و پهنای باند پورت خروجی پاسخگو نیست.
    \item \textbf{واحد واسط پورت ردیابی (\lr{TPIU - Trace Port Interface Unit}):} آخرین حلقه زنجیره درون‌تراشه‌ای که فریم‌های پروتکل \lr{ATB} را دریافت کرده، آن‌ها را با الگوهای همگام‌سازی فریم (بایت \texttt{0x7F}) بسته‌بندی نموده و به صورت جفت‌نیبل‌های موازی روی کلاک خروجی و پایه‌های خارجی تراشه منتشر می‌سازد.
\end{enumerate}

\section{پیشینه پژوهش و مقایسه تطبیقی رویکردها}
رویکردهای پایش، استخراج مسیر و ارزیابی زمان در سامانه‌های نهفته در جدول \ref{tab:trace_comparison} مورد مقایسه قرار گرفته‌اند.

\begin{table}[ht]
\centering
\small
\setlength{\tabcolsep}{3.5pt}
\renewcommand{\arraystretch}{1.45}
\caption{مقایسه تطبیقی پلتفرم‌های آزمون و ردگیری اجرای برنامه در سیستم‌های نهفته}
\label{tab:trace_comparison}
\begin{tabularx}{\textwidth}{|>{\hsize=1.3\hsize\raggedleft\arraybackslash}X|*{4}{>{\hsize=0.925\hsize\centering\arraybackslash}X|}}
\hline
\textbf{معیار مقایسه} & \textbf{ابزاردقیق نرم‌افزاری} & \textbf{شبیه‌ساز نرم‌افزاری} & \textbf{پلتفرم تجاری (\lr{Rapita/Speedgoat})} & \textbf{پلتفرم این رساله} \\
\hline
سربار زمانی \lr{(Probe Effect)} & بالا ($>15\%$) & صفر & صفر (غیرتهاجمی) & \textbf{صفر (غیرتهاجمی)} \\
\hline
تغییر کد منبع \lr{(Source Intrusion)} & نیاز دارد & بی‌نیاز & بی‌نیاز & \textbf{کاملاً بی‌نیاز} \\
\hline
دقت زمانی و رفتار چرخه‌ای & مخدوش و دارای جیتر & تقریبی & چرخه دقیق & \textbf{چرخه دقیق (جیتر صفر)} \\
\hline
هزینه لایسنس و سامانه آزمون & ناچیز & ناچیز & بسیار بالا ($>\$10,000$) & \textbf{حداقل هزینه (اقتصادی)} \\
\hline
وابستگی به برنامه کامپایل‌شده & بالا & کامل & فقط فایل باینری & \textbf{فقط باینری (\lr{.elf})} \\
\hline
پشتیبانی از پایش وقفه‌ها & نامطمئن & وابسته به مدل & کامل & \textbf{کامل و تفکیک‌شده} \\
\hline
\end{tabularx}
\end{table}

همان‌گونه که از جدول \ref{tab:trace_comparison} برمی‌آید، در صنایع هوافضا، خودروسازی و اتوماسیون پیشرفته، پلتفرم‌های تجاری آزمون خودکار محصول و تحلیل زمان‌بندی (نظیر محصولات \lr{Rapita Systems} و بسترهای \lr{Speedgoat}) استاندارد طلایی به شمار می‌روند؛ اما استقرار این پلتفرم‌ها نیازمند هزینه‌های بسیار سنگین (ده‌ها هزار دلار) و تجهیزات انحصاری است. یکی از مهم‌ترین اهداف و دستاوردهای محصول مشترک این پژوهش و پروژه همگام (سامانه آزمون خودکار)، تحقق یک پلتفرم آزمون محصول کارا و ارزان‌قیمت است که توانسته فرآیند آزمون نرم‌افزار، ماک‌سازی پریفرال‌ها و استخراج مسیر اجرا را با حداقل هزینه محقق سازد.

همچنین در سطح سخت‌افزار دریافت سیگنال، اگرچه پروب‌های سخت‌افزاری تجاری نظیر \lr{SEGGER J-Trace} به ثبت مستقیم در فرکانس‌های بالا می‌پردازند، اما درون آن‌ها از تراشه‌های اختصاصی \lr{FPGA} و حافظه‌های سریع \lr{DRAM} استفاده شده است. آزمون‌های تجربی این پروژه نشان داد که لاجیک‌آنالایزرهای عمومی اقتصادی نهایتاً برای فرکانس‌های بسیار بالا ایده‌آل نبوده و بهترین انتخاب سخت‌افزاری، یک برد سفارشی مبتنی بر \lr{FPGA} است؛ هرچند هزینه پیاده‌سازی اختصاصی آن با هزینه خرید پروب‌های تجاری تقریباً برابری خواهد کرد. با این حال، ارزش افزوده اصلی این رساله، توسعه کل چرخه نرم‌افزاری استخراج و دیکودینگ داده‌ها با هزینه بسیار پایین است.

\section{جمع‌بندی فصل}
در این فصل پایه‌های تئوریک محاسبه \lr{WCET} و ساختار لایه‌ای معماری \lr{ARM CoreSight} تشریح گردید. همچنین ضرورت اجتناب از ابزاردقیق نرم‌افزاری و مزایای اقتصادی و فنی پلتفرم توسعه‌داده‌شده نسبت به راهکارهای تجاری گران‌قیمت بازار مورد ارزیابی قرار گرفت. با استناد به این مبانی، در فصل سوم به طراحی معماری کلان سامانه و پیاده‌سازی گام‌به‌گام خط لوله تأمین داده پرداخته خواهد شد.
